\documentclass[11pt,a4paper]{article}

\usepackage[utf8]{inputenc}
\usepackage[T1]{fontenc}
\usepackage[french]{babel}
\usepackage[a4paper,margin=2cm]{geometry}
\usepackage{lmodern}
\usepackage{amsmath,amssymb,amsthm,mathtools}
\usepackage{enumitem}
\usepackage{fancyhdr}
\usepackage{lastpage}
\usepackage{xcolor}
\usepackage{tikz}
\usepackage{tcolorbox}
\usepackage{hyperref}
\usepackage{stmaryrd}
\usepackage{array}
\usepackage{mathrsfs}
\usepackage{multicol}
\usepackage{titlesec}

\hypersetup{colorlinks=true,linkcolor=blue!50!black,urlcolor=blue!50!black}

\definecolor{bleu}{RGB}{35,72,130}
\definecolor{gris}{RGB}{80,80,80}
\definecolor{vertcorr}{RGB}{20,110,70}
\definecolor{fondclair}{RGB}{245,247,250}

\setlength{\headheight}{14pt}
\pagestyle{fancy}
\fancyhf{}
\lhead{Algèbre linéaire --DM1 pgb}
\rhead{Sujet + corrigé }
\cfoot{\thepage/\pageref{LastPage}}

\titleformat{\section}{\Large\bfseries\color{bleu}}{}{0pt}{}
\titleformat{\subsection}{\large\bfseries\color{bleu}}{}{0pt}{}
\titleformat{\subsubsection}{\normalsize\bfseries\color{gris}}{}{0pt}{}

\newtheorem*{remarque}{Remarque}
\newtheorem*{methode}{Méthode}

\newtcolorbox{rappelbox}{
  colback=fondclair,
  colframe=bleu,
  boxrule=0.8pt,
  arc=2mm,
  left=2mm,right=2mm,top=1.5mm,bottom=1.5mm
}

\newtcolorbox{corrbox}{
  colback=green!2,
  colframe=vertcorr,
  boxrule=0.8pt,
  arc=2mm,
  left=2mm,right=2mm,top=1.5mm,bottom=1.5mm
}

\newcommand{\R}{\mathbb{R}}
\newcommand{\Vect}{\operatorname{Vect}}
\newcommand{\dimv}{\operatorname{dim}}
\newcommand{\ds}{\displaystyle}
\newcommand{\cours}[1]{\textit{[Réf. cours : #1]}}
\newcommand{\p}[1]{\medskip\noindent\textbf{#1.}}

\begin{document}

\begin{center}
    {\LARGE\bfseries\color{bleu} Sujet d'algèbre linéaire -- niveau PCSI}\\[0.4em]
    {\large Deux problèmes classiques, un problème plus élégant}\\[0.6em]
    \rule{0.85\linewidth}{0.8pt}
\end{center}

\vspace{0.4em}

\begin{rappelbox}
\textbf{Rappels de notions utiles, pas forcément développées dans les chapitres fournis.}
\begin{enumerate}[label=\arabic*)]
    \item \textbf{Plan affine de l'espace.} Si $A$ est un point de l'espace et si $u,v$ sont deux vecteurs non colinéaires, alors l'ensemble des points $M$ tels que
    \[
    \overrightarrow{AM}=\lambda u+\mu v \qquad (\lambda,\mu\in\R)
    \]
    est un plan affine.
    
    \item \textbf{Barycentre de trois points de mêmes coefficients.} Pour des points $P,Q,R$, le barycentre affecté des coefficients $(1,1,1)$ est le point $G$ vérifiant
    \[
    \overrightarrow{OG}=\frac13\bigl(\overrightarrow{OP}+\overrightarrow{OQ}+\overrightarrow{OR}\bigr)
    \]
    pour tout origine $O$. On a alors
    \[
    \overrightarrow{GP}+\overrightarrow{GQ}+\overrightarrow{GR}=0.
    \]
    
    \item \textbf{Centre de symétrie.} Un point $G$ est centre de symétrie de deux points $P$ et $P'$ si et seulement si
    \[
    \overrightarrow{GP'}=-\overrightarrow{GP}.
    \]
    Cela équivaut à dire que $G$ est le milieu de $[PP']$.
    
    \item \textbf{Produit scalaire et norme dans $\R^3$.} Pour $a,b\in\R^3$,
    \[
    \|a-b\|^2=(a-b)\cdot(a-b)=\|a\|^2+\|b\|^2-2a\cdot b.
    \]
    Si deux vecteurs non nuls $x$ et $y$ vérifient
    \[
    x\cdot y=\|x\|\,\|y\|\cos\theta,
    \]
    alors $\theta$ est l'angle entre eux.
    
    \item \textbf{Déterminant dans le plan.} Dans une base du plan, deux vecteurs sont libres si et seulement si le déterminant de leurs coordonnées dans cette base est non nul.
\end{enumerate}
\end{rappelbox}

\vspace{0.6em}

\begin{center}
\textbf{Énoncé}
\end{center}

\section*{Problème 1 -- Sous-espaces de $\R^4$, intersections, sommes et bases}

Dans tout le problème, on travaille dans $E=\R^4$.

On considère les sous-espaces
\[
F=\{(x,y,z,t)\in\R^4\mid x+y+z+t=0\text{ et }x-z=0\}
\]
et
\[
G=\Vect\bigl((1,1,1,1),(0,1,0,-1)\bigr).
\]

\subsection*{Partie I -- Étude de $F$}
\begin{enumerate}[label=\arabic*.]
    \item Montrer que $F$ est un sous-espace vectoriel de $\R^4$.
    \item Déterminer une écriture paramétrique des éléments de $F$.
    \item En déduire une famille génératrice de $F$.
    \item Montrer que cette famille est libre.
    \item Donner une base de $F$ et déterminer $\dimv(F)$.
\end{enumerate}

\subsection*{Partie II -- Étude de $G$}
\begin{enumerate}[label=\arabic*.]
    \setcounter{enumi}{5}
    \item Montrer que $G$ est un sous-espace vectoriel de $\R^4$.
    \item Déterminer une écriture paramétrique des éléments de $G$.
    \item Montrer que la famille donnée dans la définition de $G$ est libre.
    \item En déduire une base de $G$ et sa dimension.
\end{enumerate}

\subsection*{Partie III -- Intersection et somme}
\begin{enumerate}[label=\arabic*.]
    \setcounter{enumi}{9}
    \item Déterminer $F\cap G$.
    \item En donner une base et sa dimension.
    \item Les sous-espaces $F$ et $G$ sont-ils en somme directe ?
    \item Calculer $\dimv(F+G)$.
    \item A-t-on $F+G=\R^4$ ?
\end{enumerate}

\subsection*{Partie IV -- Une famille de vecteurs}
On pose
\[
u_1=(1,0,1,-2),\qquad u_2=(0,1,0,-1),\qquad u_3=(1,1,1,1).
\]
\begin{enumerate}[label=\arabic*.]
    \setcounter{enumi}{14}
    \item Montrer que $(u_1,u_2)$ est une base de $F$.
    \item Montrer que $(u_2,u_3)$ est une base de $G$.
    \item Étudier la famille $(u_1,u_2,u_3)$ : est-elle libre ?
    \item Que peut-on en déduire sur $F+G$ ?
    \item Extraire, à partir de $(u_1,u_2,u_3)$, une base de $F+G$.
\end{enumerate}

\subsection*{Partie V -- Décomposition d'un vecteur}
\begin{enumerate}[label=\arabic*.]
    \setcounter{enumi}{19}
    \item Soit $w=(x,y,z,t)\in\R^4$. Montrer que si $w\in F+G$, alors on peut écrire
    \[
    w=\alpha u_1+\beta u_2+\gamma u_3.
    \]
    
    \item En identifiant les coordonnées, montrer que tout vecteur de $F+G$ vérifie une relation linéaire simple entre ses coordonnées.
    
    \item Réciproquement, montrer que tout vecteur de $\R^4$ vérifiant cette relation appartient à $F+G$.
    
    \item En déduire un système d'équations cartésiennes de $F+G$.
\end{enumerate}

\section*{Problème 2 -- Une figure hexagonale dans un plan de l'espace}

Dans l'espace vectoriel $\R^3$, on considère trois vecteurs $u,v,w$ formant une famille libre.
On définit six points par leurs vecteurs-position :
\[
A=u,\qquad B=v,\qquad C=w,\qquad D=v+w-u,\qquad E=w+u-v,\qquad F=u+v-w.
\]

On s'intéresse à la figure $ABCDEF$.

\subsection*{Partie I -- Calculs de vecteurs et premières propriétés}
\begin{enumerate}[label=\arabic*.]
    \item Exprimer les vecteurs
    \[
    \overrightarrow{AB},\ \overrightarrow{BC},\ \overrightarrow{CD},\ \overrightarrow{DE},\ \overrightarrow{EF},\ \overrightarrow{FA}
    \]
    en fonction de $u,v,w$.
    \item Montrer que
    \[
    \overrightarrow{AB}=\overrightarrow{CD},\qquad \overrightarrow{BC}=\overrightarrow{FA}.
    \]
    \item Montrer que
    \[
    \overrightarrow{DE}=-2\overrightarrow{AB},\qquad \overrightarrow{EF}=-2\overrightarrow{BC}.
    \]
    \item Montrer que
    \[
    \overrightarrow{AB}+\overrightarrow{BC}+\overrightarrow{CA}=0.
    \]
    \item Interpréter géométriquement cette relation.
\end{enumerate}

\subsection*{Partie II -- Le plan de la figure}
\begin{enumerate}[label=\arabic*.]
    \setcounter{enumi}{5}
    \item Montrer que les six points $A,B,C,D,E,F$ appartiennent à un même plan affine.
    \item Donner une représentation paramétrique de ce plan à partir du point $A$.
    \item Montrer que $A,B,C$ ne sont pas alignés.
    \item En déduire que ce plan est exactement le plan affine engendré par $A,B,C$.
\end{enumerate}

\subsection*{Partie III -- Deux centres remarquables}
On pose
\[
G=\frac13(u+v+w),\qquad S=\frac12(u+v+w).
\]
\begin{enumerate}[label=\arabic*.]
    \setcounter{enumi}{9}
    \item Montrer que $G$ est le barycentre des points $A,B,C$ affectés de coefficients égaux.
    \item Montrer aussi que $G$ est le barycentre des points $D,E,F$ affectés de coefficients égaux.
    \item Montrer que $S$ est le milieu de $[AD]$, puis de $[BE]$ et de $[CF]$.
    \item En déduire que la symétrie centrale de centre $S$ échange $A$ et $D$, $B$ et $E$, $C$ et $F$.
    \item Conclure que la figure $ABCDEF$ admet un centre de symétrie.
\end{enumerate}

\subsection*{Partie IV -- Coordonnées affines dans le plan}
On considère le repère affine du plan $\mathcal P=(A;\overrightarrow{AB},\overrightarrow{AC})$.
\begin{enumerate}[label=\arabic*.]
    \setcounter{enumi}{14}
    \item Justifier que tout point $M$ du plan $\mathcal P$ possède des coordonnées affines uniques $(\alpha,\beta)$ définies par
    \[
    \overrightarrow{AM}=\alpha\,\overrightarrow{AB}+\beta\,\overrightarrow{AC}.
    \]
    \item Déterminer les coordonnées affines des points $A,B,C,D,E,F$ dans ce repère.
    \item Montrer que $ABCD$ est un parallélogramme.
    \item Déterminer les coordonnées affines de $G$ et de $S$ dans ce repère.
    \item Retrouver, à partir de ces coordonnées, que $S$ est centre de symétrie de la figure.
\end{enumerate}

\subsection*{Partie V -- Une écriture vectorielle élégante}
\begin{enumerate}[label=\arabic*.]
    \setcounter{enumi}{19}
    \item Montrer que les vecteurs $u-v$ et $v-w$ sont libres.
    \item Montrer que tout vecteur directeur du plan de la figure s'écrit de manière unique sous la forme
    \[
    \alpha(u-v)+\beta(v-w),\qquad (\alpha,\beta)\in\R^2.
    \]
    \item Exprimer $\overrightarrow{SA}$, $\overrightarrow{SB}$, $\overrightarrow{SC}$, $\overrightarrow{SD}$, $\overrightarrow{SE}$ et $\overrightarrow{SF}$ en fonction de $u-v$ et $v-w$.
    \item Montrer que les six points se regroupent par paires de vecteurs opposés par rapport à $S$.
\end{enumerate}

\subsection*{Partie VI -- Une hypothèse métrique naturelle}
On suppose dans cette partie que
\[
\|u-v\|=\|v-w\|=\|w-u\|.
\]
\begin{enumerate}[label=\arabic*.]
    \setcounter{enumi}{23}
    \item Montrer que le triangle $ABC$ est équilatéral.
    \item Montrer que le triangle $DEF$ est équilatéral.
    \item Montrer que
    \[
    AB=BC=CD=FA\qquad\text{et}\qquad DE=EF=FD=2AB.
    \]
    \item Interpréter géométriquement la figure obtenue.
\end{enumerate}


\section*{Problème 3 -- Élégant : quand deux sous-espaces engendrés coïncident}

Dans un espace vectoriel réel $E$, on considère trois vecteurs $u,v,w$.
On suppose que la famille $(u,v)$ est libre.
On pose
\[
F=\Vect(u+v,\,v+w),\qquad G=\Vect(u,v,w).
\]

L'objectif est de déterminer dans quels cas on a $F=G$.

\subsection*{Partie I -- Premières observations}
\begin{enumerate}[label=\arabic*.]
    \item Montrer que $F\subset G$.
    \item Montrer que $\dimv(G)\in\{2,3\}$.
    \item Donner une condition nécessaire et suffisante pour que $\dimv(G)=2$.
\end{enumerate}

\subsection*{Partie II -- Cas où $w\in\Vect(u,v)$}
On suppose dans cette partie qu'il existe $(a,b)\in\R^2$ tel que
\[
w=au+bv.
\]
\begin{enumerate}[label=\arabic*.]
    \setcounter{enumi}{3}
    \item Exprimer $v+w$ en fonction de $u$ et $v$.
    \item Écrire les coordonnées de $u+v$ et $v+w$ dans la base $(u,v)$.
    \item Montrer que $F=\Vect(u,v)$ si et seulement si
    \[
    \det\begin{pmatrix}
    1 & a\\
    1 & 1+b
    \end{pmatrix}\neq 0.
    \]
    \item Simplifier cette condition.
    \item En déduire une condition nécessaire et suffisante sur $a$ et $b$ pour que $F=G$.
\end{enumerate}

\subsection*{Partie III -- Cas où $(u,v,w)$ est libre}
\begin{enumerate}[label=\arabic*.]
    \setcounter{enumi}{8}
    \item Montrer que $\dimv(G)=3$.
    \item Montrer que $\dimv(F)\leq 2$.
    \item En déduire que $F\neq G$.
\end{enumerate}

\subsection*{Partie IV -- Conclusion générale}
\begin{enumerate}[label=\arabic*.]
    \setcounter{enumi}{11}
    \item Montrer que
    \[
    F=G \iff \bigl(w\in\Vect(u,v)\text{ et }u+v,\,v+w\text{ forment une base de }\Vect(u,v)\bigr).
    \]
    \item Donner une condition finale simple sous la forme
    \[
    F=G\iff w=au+bv\text{ avec }\cdots
    \]
\end{enumerate}

\subsection*{Partie V -- Retour à $\R^2$}
On se place dans $E=\R^2$ avec
\[
u=(1,0),\qquad v=(0,1),\qquad w=(a,b).
\]
\begin{enumerate}[label=\arabic*.]
    \setcounter{enumi}{13}
    \item Déterminer explicitement une condition sur $a$ et $b$ pour que
    \[
    \Vect(u+v,\,v+w)=\R^2.
    \]
    \item Retrouver cette condition par un calcul direct de déterminant dans la base canonique.
    \item Donner un exemple où $F=G$, puis un exemple où $F\neq G$.
\end{enumerate}

\newpage
\begin{center}
\textbf{Corrigé détaillé}
\end{center}

\section*{Corrigé du problème 1}


\p{1} Montrons que $F$ est un sous-espace vectoriel de $\R^4$.

On peut utiliser le critère de sous-espace vectoriel \cours{chap. 3, th. 5.3}.
L'ensemble $F$ n'est pas vide car $(0,0,0,0)\in F$.

Soient $X=(x,y,z,t)$ et $X'=(x',y',z',t')$ dans $F$, et soit $\lambda\in\R$.
Alors
\[
 x+y+z+t=0,\qquad x-z=0,
\]
et
\[
 x'+y'+z'+t'=0,\qquad x'-z'=0.
\]
Pour la somme,
\[
(x+x')+(y+y')+(z+z')+(t+t')=(x+y+z+t)+(x'+y'+z'+t')=0,
\]
et
\[
(x+x')-(z+z')=(x-z)+(x'-z')=0.
\]
Donc $X+X'\in F$.

Pour le produit par un scalaire,
\[
\lambda x+\lambda y+\lambda z+\lambda t=\lambda(x+y+z+t)=0,
\qquad
\lambda x-\lambda z=\lambda(x-z)=0.
\]
Donc $\lambda X\in F$.

Ainsi $F$ est un sous-espace vectoriel de $\R^4$.

\p{2} Cherchons une écriture paramétrique.

La condition $x-z=0$ donne
\[
z=x.
\]
En remplaçant dans $x+y+z+t=0$, on obtient
\[
2x+y+t=0,
\qquad\text{donc}\qquad t=-2x-y.
\]
En posant $x=a$ et $y=b$, on obtient
\[
(x,y,z,t)=(a,b,a,-2a-b).
\]
Donc
\[
F=\{(a,b,a,-2a-b)\mid (a,b)\in\R^2\}.
\]

\p{3} On réécrit
\[
(a,b,a,-2a-b)=a(1,0,1,-2)+b(0,1,0,-1).
\]
Donc
\[
F=\Vect\bigl((1,0,1,-2),(0,1,0,-1)\bigr).
\]
La famille
\[
\bigl((1,0,1,-2),(0,1,0,-1)\bigr)
\]
est donc génératrice de $F$ \cours{chap. 4, déf. 3.1}.

\p{4} Montrons qu'elle est libre.

Supposons
\[
\alpha(1,0,1,-2)+\beta(0,1,0,-1)=(0,0,0,0).
\]
En identifiant les coordonnées, on obtient
\[
\alpha=0,\qquad \beta=0.
\]
La seule combinaison linéaire nulle est donc la combinaison triviale ; la famille est libre \cours{chap. 4, déf. 4.1}.

\p{5} Une famille libre et génératrice de $F$ est une base de $F$ \cours{chap. 4, déf. 7.1}. Ainsi,
\[
\mathcal B_F=\bigl((1,0,1,-2),(0,1,0,-1)\bigr)
\]
est une base de $F$, et
\[
\dimv(F)=2.
\]

\p{6} Par définition,
\[
G=\Vect\bigl((1,1,1,1),(0,1,0,-1)\bigr).
\]
Tout sous-espace engendré est un sous-espace vectoriel \cours{chap. 3, §7.2}. Donc $G$ est un sous-espace vectoriel de $\R^4$.

\p{7} Tout vecteur de $G$ s'écrit
\[
\alpha(1,1,1,1)+\beta(0,1,0,-1)=(\alpha,\alpha+\beta,\alpha,\alpha-\beta).
\]
Donc
\[
G=\{(\alpha,\alpha+\beta,\alpha,\alpha-\beta)\mid (\alpha,\beta)\in\R^2\}.
\]

\p{8} Supposons
\[
\alpha(1,1,1,1)+\beta(0,1,0,-1)=0.
\]
En regardant la première coordonnée, on obtient $\alpha=0$, puis la deuxième donne $\beta=0$.
La famille est donc libre.

\p{9} La famille donnée est libre et génératrice de $G$, donc c'est une base de $G$, et
\[
\dimv(G)=2.
\]

\p{10} Déterminons $F\cap G$.

Soit $X\in F\cap G$. Comme $X\in G$, il existe $\alpha,\beta\in\R$ tels que
\[
X=(\alpha,\alpha+\beta,\alpha,\alpha-\beta).
\]
La condition d'appartenance à $F$ donne
\[
x+y+z+t=0.
\]
On calcule :
\[
\alpha+(\alpha+\beta)+\alpha+(\alpha-\beta)=4\alpha.
\]
Donc $4\alpha=0$, soit $\alpha=0$.
Alors
\[
X=(0,\beta,0,-\beta)=\beta(0,1,0,-1).
\]
Réciproquement, tout vecteur de la forme $\beta(0,1,0,-1)$ appartient manifestement à $F$ et à $G$.
Donc
\[
F\cap G=\Vect(0,1,0,-1).
\]

\p{11} Une base de $F\cap G$ est
\[
\bigl((0,1,0,-1)\bigr),
\]
donc
\[
\dimv(F\cap G)=1.
\]

\p{12} Deux sous-espaces sont en somme directe si et seulement si leur intersection est réduite à $\{0\}$ \cours{chap. 5 dimension finie, déf. 10.1}. Ici,
\[
F\cap G=\Vect(0,1,0,-1)\neq\{0\}.
\]
Donc $F$ et $G$ ne sont pas en somme directe.

\p{13} Par Grassmann \cours{chap. 5 dimension finie, th. 9.1},
\[
\dimv(F+G)=\dimv(F)+\dimv(G)-\dimv(F\cap G)=2+2-1=3.
\]

\p{14} Comme $\dimv(\R^4)=4$ et que $\dimv(F+G)=3$, on n'a pas
\[
F+G=\R^4.
\]

\p{15} D'après les questions précédentes, $(u_1,u_2)$ est exactement la famille génératrice libre trouvée pour $F$ ; c'est donc une base de $F$.

\p{16} Le vecteur $u_2=(0,1,0,-1)$ et le vecteur $u_3=(1,1,1,1)$ sont précisément les deux générateurs de $G$ ; ils forment donc une base de $G$.

\p{17} Étudions la liberté de $(u_1,u_2,u_3)$.

Supposons
\[
\alpha u_1+\beta u_2+\gamma u_3=0.
\]
Cela donne
\[
\alpha(1,0,1,-2)+\beta(0,1,0,-1)+\gamma(1,1,1,1)=(0,0,0,0).
\]
En identifiant la première coordonnée :
\[
\alpha+\gamma=0.
\]
La deuxième donne
\[
\beta+\gamma=0.
\]
La quatrième donne
\[
-2\alpha-\beta+\gamma=0.
\]
En remplaçant $\gamma=-\alpha$ et $\beta=\alpha$, on obtient
\[
-2\alpha-\alpha-\alpha=-4\alpha=0,
\]
donc $\alpha=0$, puis $\beta=\gamma=0$.
La famille $(u_1,u_2,u_3)$ est donc libre.

\p{18} Cette famille est constituée de vecteurs appartenant à $F+G$ et contient $3$ vecteurs libres. Or on a vu que
\[
\dimv(F+G)=3.
\]
Donc $(u_1,u_2,u_3)$ est une base de $F+G$ \cours{chap. 5 dimension finie, §12.3}.

\p{19} Une base extraite de $(u_1,u_2,u_3)$ de $F+G$ est donc simplement
\[
\bigl(u_1,u_2,u_3\bigr).
\]

\p{20} Si $w\in F+G$, alors par définition de la somme de sous-espaces,
\[
w=f+g,
\qquad f\in F,\ g\in G.
\]
Or $(u_1,u_2)$ est une base de $F$ et $(u_2,u_3)$ une base de $G$. On peut donc écrire
\[
f=au_1+bu_2,
\qquad g=cu_2+du_3.
\]
Ainsi
\[
w=au_1+(b+c)u_2+du_3.
\]
En posant $\alpha=a$, $\beta=b+c$, $\gamma=d$, on obtient bien
\[
w=\alpha u_1+\beta u_2+\gamma u_3.
\]

\p{21} Écrivons
\[
(x,y,z,t)=\alpha(1,0,1,-2)+\beta(0,1,0,-1)+\gamma(1,1,1,1).
\]
Alors
\[
(x,y,z,t)=(\alpha+\gamma,\beta+\gamma,\alpha+\gamma,-2\alpha-\beta+\gamma).
\]
On voit immédiatement que
\[
x=z.
\]
Tout vecteur de $F+G$ vérifie donc la relation linéaire
\[
x-z=0.
\]

\p{22} Réciproquement, soit $(x,y,z,t)\in\R^4$ tel que $x=z$.
Posons
\[
\gamma=x,
\qquad \beta=y-x,
\qquad \alpha=0.
\]
Alors
\[
\alpha+\gamma=x,
\qquad \beta+\gamma=y,
\qquad \alpha+\gamma=x=z,
\]
et la quatrième coordonnée vaut
\[
-2\alpha-\beta+\gamma=-(y-x)+x=2x-y.
\]
On obtient donc exactement $(x,y,z,t)$ si et seulement si $t=2x-y$.
Mais comme $x=z$, cette relation équivaut à
\[
x-y+z-t=0.
\]
On peut aussi procéder directement : en utilisant la base $(u_1,u_2,u_3)$, tout élément de $F+G$ est un élément d'un espace de dimension $3$ défini par une seule équation linéaire ; ici la bonne relation est
\[
x-z=0.
\]
Comme $\dimv(F+G)=3$, cette seule relation suffit à décrire $F+G$.
En effet, l'ensemble
\[
H=\{(x,y,z,t)\in\R^4\mid x-z=0\}
\]
est un hyperplan de $\R^4$, donc de dimension $3$, et contient $u_1,u_2,u_3$. Il contient donc $F+G$, et les dimensions étant égales, on a $F+G=H$.

\p{23} Finalement,
\[
F+G=\{(x,y,z,t)\in\R^4\mid x-z=0\}.
\]
Il s'agit donc d'un hyperplan de $\R^4$.

\section*{Corrigé du problème 2}

\p{1} On calcule directement :
\[
\overrightarrow{AB}=v-u,
\qquad
\overrightarrow{BC}=w-v,
\qquad
\overrightarrow{CD}=D-C=(v+w-u)-w=v-u,
\]
\[
\overrightarrow{DE}=E-D=(w+u-v)-(v+w-u)=2u-2v,
\]
\[
\overrightarrow{EF}=F-E=(u+v-w)-(w+u-v)=2v-2w,
\qquad
\overrightarrow{FA}=A-F=u-(u+v-w)=w-v.
\]

\p{2} D'après la question précédente,
\[
\overrightarrow{AB}=v-u=\overrightarrow{CD},
\qquad
\overrightarrow{BC}=w-v=\overrightarrow{FA}.
\]

\p{3} On reprend les calculs de la question 1 :
\[
\overrightarrow{DE}=2u-2v=-2(v-u)=-2\overrightarrow{AB},
\]
\[
\overrightarrow{EF}=2v-2w=-2(w-v)=-2\overrightarrow{BC}.
\]

\p{4} La relation de Chasles donne
\[
\overrightarrow{AB}+\overrightarrow{BC}=\overrightarrow{AC}.
\]
Donc
\[
\overrightarrow{AB}+\overrightarrow{BC}+\overrightarrow{CA}=0.
\]
En remplaçant par les expressions trouvées, on obtient aussi
\[
(v-u)+(w-v)+(u-w)=0.
\]

\p{5} Cette relation signifie que les vecteurs $\overrightarrow{AB}$, $\overrightarrow{BC}$ et $\overrightarrow{CA}$ forment le contour fermé du triangle $ABC$.

\p{6} Soit
\[
\mathcal P=A+\Vect(\overrightarrow{AB},\overrightarrow{AC})=u+\Vect(v-u,w-u).
\]
Montrons que $D,E,F\in\mathcal P$.
On a
\[
\overrightarrow{AD}=D-A=(v+w-u)-u=(v-u)+(w-u),
\]
\[
\overrightarrow{AE}=E-A=(w+u-v)-u=w-v=(w-u)-(v-u),
\]
\[
\overrightarrow{AF}=F-A=(u+v-w)-u=v-w=(v-u)-(w-u).
\]
Ainsi, $\overrightarrow{AD}$, $\overrightarrow{AE}$ et $\overrightarrow{AF}$ appartiennent à $\Vect(v-u,w-u)$, d'où $D,E,F\in\mathcal P$.
Les six points sont donc coplanaires.

\p{7} Une représentation paramétrique de ce plan est
\[
\mathcal P=\{u+\lambda(v-u)+\mu(w-u)\mid (\lambda,\mu)\in\R^2\}.
\]

\p{8} Si $A,B,C$ étaient alignés, les vecteurs $v-u$ et $w-u$ seraient colinéaires. Il existerait donc $(\lambda,\mu)\neq(0,0)$ tel que
\[
\lambda(v-u)+\mu(w-u)=0.
\]
Cela s'écrit
\[
(-\lambda-\mu)u+\lambda v+\mu w=0.
\]
Comme la famille $(u,v,w)$ est libre, on obtient
\[
-\lambda-\mu=0,\qquad \lambda=0,\qquad \mu=0,
\]
contradiction. Donc $A,B,C$ ne sont pas alignés.

\p{9} Trois points non alignés déterminent un unique plan affine. Le plan trouvé à la question 7 est donc exactement le plan affine engendré par $A,B,C$.

\p{10} On a, par définition des points,
\[
A+B+C=u+v+w.
\]
Donc
\[
G=\frac13(u+v+w)=\frac13(A+B+C).
\]
D'après le rappel de début de sujet, $G$ est bien le barycentre de $A,B,C$ affectés des coefficients $(1,1,1)$.

\p{11} On additionne maintenant les vecteurs-position de $D,E,F$ :
\[
D+E+F=(v+w-u)+(w+u-v)+(u+v-w)=u+v+w.
\]
Ainsi,
\[
G=\frac13(D+E+F).
\]
Le point $G$ est donc aussi le barycentre de $D,E,F$ affectés de coefficients égaux.

\p{12} On calcule :
\[
A+D=u+(v+w-u)=u+v+w=2S.
\]
Donc $S$ est le milieu de $[AD]$.
De même,
\[
B+E=v+(w+u-v)=u+v+w=2S,
\]
\[
C+F=w+(u+v-w)=u+v+w=2S.
\]
Donc $S$ est aussi le milieu de $[BE]$ et de $[CF]$.

\p{13} Si $S$ est le milieu de $[AD]$, la symétrie centrale de centre $S$ échange $A$ et $D$. De même, elle échange $B$ et $E$, puis $C$ et $F$.

\p{14} La figure $ABCDEF$ admet donc un centre de symétrie, qui est
\[
S=\frac12(u+v+w).
\]

\p{15} D'après la question 8, les vecteurs $\overrightarrow{AB}$ et $\overrightarrow{AC}$ ne sont pas colinéaires. Ils forment donc une base de l'espace vectoriel directeur du plan $\mathcal P$. Par conséquent, tout point $M\in\mathcal P$ s'écrit de manière unique sous la forme
\[
\overrightarrow{AM}=\alpha\overrightarrow{AB}+\beta\overrightarrow{AC}.
\]

\p{16} On obtient immédiatement
\[
A(0,0),\qquad B(1,0),\qquad C(0,1).
\]
De plus,
\[
\overrightarrow{AD}=(v-u)+(w-u)=\overrightarrow{AB}+\overrightarrow{AC},
\]
donc
\[
D(1,1).
\]
Ensuite,
\[
\overrightarrow{AE}=w-v=(w-u)-(v-u)=\overrightarrow{AC}-\overrightarrow{AB},
\]
donc
\[
E(-1,1).
\]
Enfin,
\[
\overrightarrow{AF}=v-w=(v-u)-(w-u)=\overrightarrow{AB}-\overrightarrow{AC},
\]
donc
\[
F(1,-1).
\]

\p{17} Dans ce repère affine,
\[
A(0,0),\quad B(1,0),\quad C(0,1),\quad D(1,1).
\]
On voit que
\[
\overrightarrow{AB}=(1,0)=\overrightarrow{CD},
\qquad
\overrightarrow{AC}=(0,1)=\overrightarrow{BD}.
\]
Le quadrilatère $ABCD$ est donc un parallélogramme.

\p{18} Comme
\[
G=\frac13(A+B+C),
\]
on lit directement dans le repère affine
\[
G\left(\frac13,\frac13\right).
\]
De même,
\[
S=\frac12(A+D),
\]
donc
\[
S\left(\frac12,\frac12\right).
\]

\p{19} Dans ce repère affine, on vérifie facilement que
\[
\frac{A+D}{2}=\frac{(0,0)+(1,1)}{2}=\left(\frac12,\frac12\right)=S,
\]
\[
\frac{B+E}{2}=\frac{(1,0)+(-1,1)}{2}=\left(\frac12,\frac12\right)=S,
\]
\[
\frac{C+F}{2}=\frac{(0,1)+(1,-1)}{2}=\left(\frac12,\frac12\right)=S.
\]
On retrouve ainsi que $S$ est le milieu commun de $[AD]$, $[BE]$ et $[CF]$, donc le centre de symétrie de la figure.

\p{20} Supposons
\[
\alpha(u-v)+\beta(v-w)=0.
\]
Alors
\[
\alpha u+(-\alpha+\beta)v-\beta w=0.
\]
Comme $(u,v,w)$ est libre, tous les coefficients sont nuls, donc
\[
\alpha=0,\qquad \beta=0.
\]
Les vecteurs $u-v$ et $v-w$ sont donc libres.

\p{21} Le plan directeur de $\mathcal P$ est engendré par $v-u$ et $w-u$, donc aussi par $u-v$ et $v-w$. D'après la question 20, ces deux vecteurs sont libres ; ils forment donc une base du plan directeur. Tout vecteur directeur du plan s'écrit donc de manière unique sous la forme
\[
\alpha(u-v)+\beta(v-w).
\]

\p{22} Posons
\[
a=u-v,
\qquad
b=v-w.
\]
Alors $u-w=a+b$.
Comme $S=\frac12(u+v+w)$, on calcule
\[
\overrightarrow{SA}=u-S=\frac12(u-v)+\frac12(u-w)=a+\frac12 b,
\]
\[
\overrightarrow{SB}=v-S=\frac12(v-u)+\frac12(v-w)=-\frac12 a+\frac12 b,
\]
\[
\overrightarrow{SC}=w-S=\frac12(w-u)+\frac12(w-v)=-\frac12 a-b.
\]
Comme $S$ est centre de symétrie, on en déduit aussitôt
\[
\overrightarrow{SD}=-\overrightarrow{SA}=-a-\frac12 b,
\]
\[
\overrightarrow{SE}=-\overrightarrow{SB}=\frac12 a-\frac12 b,
\]
\[
\overrightarrow{SF}=-\overrightarrow{SC}=\frac12 a+b.
\]

\p{23} La question précédente donne immédiatement
\[
\overrightarrow{SD}=-\overrightarrow{SA},
\qquad
\overrightarrow{SE}=-\overrightarrow{SB},
\qquad
\overrightarrow{SF}=-\overrightarrow{SC}.
\]
Les six points se regroupent donc par paires de vecteurs opposés par rapport à $S$ : $(A,D)$, $(B,E)$ et $(C,F)$.

\p{24} L'hypothèse
\[
\|u-v\|=\|v-w\|=\|w-u\|
\]
se traduit exactement par
\[
AB=BC=CA.
\]
Le triangle $ABC$ est donc équilatéral.

\p{25} Calculons les côtés du triangle $DEF$ :
\[
\overrightarrow{DE}=2(u-v),
\qquad
\overrightarrow{EF}=2(v-w),
\qquad
\overrightarrow{FD}=D-F=(v+w-u)-(u+v-w)=2(w-u).
\]
Donc
\[
DE=2\|u-v\|,
\qquad
EF=2\|v-w\|,
\qquad
FD=2\|w-u\|.
\]
Sous l'hypothèse faite, ces trois longueurs sont égales. Le triangle $DEF$ est donc lui aussi équilatéral.

\p{26} On a déjà
\[
AB=\|u-v\|,
\qquad
BC=\|v-w\|,
\qquad
CD=\|v-u\|=AB,
\qquad
FA=\|w-v\|=BC.
\]
Sous l'hypothèse métrique, cela donne
\[
AB=BC=CD=FA.
\]
De plus, d'après la question 25,
\[
DE=EF=FD=2AB.
\]
Ainsi,
\[
AB=BC=CD=FA
\qquad\text{et}\qquad
DE=EF=FD=2AB.
\]

\p{27} Géométriquement, la figure est constituée de deux triangles équilatéraux coplanaires : le triangle $ABC$ et le triangle $DEF$. Ils ont le même barycentre $G$ d'après les questions 10 et 11, et ils sont échangés par la symétrie centrale de centre $S$ d'après les questions 12 à 14.


\section*{Corrigé du problème 3}


\p{1} Comme $u+v$ et $v+w$ sont chacun des combinaisons linéaires de $u,v,w$, on a
\[
u+v\in G,
\qquad
v+w\in G.
\]
Donc tout vecteur de $F=\Vect(u+v,v+w)$ appartient à $G$. Ainsi
\[
F\subset G.
\]

\p{2} Puisque $(u,v)$ est libre, on a
\[
\dimv\Vect(u,v)=2.
\]
Or $G=\Vect(u,v,w)$ contient $u$ et $v$ ; sa dimension est donc au moins $2$. Mais il est engendré par trois vecteurs, donc sa dimension est au plus $3$.
Ainsi
\[
\dimv(G)\in\{2,3\}.
\]

\p{3} On a
\[
\dimv(G)=2
\iff w\in\Vect(u,v).
\]
En effet, si $w\in\Vect(u,v)$, alors $G=\Vect(u,v)$, donc $\dimv(G)=2$. Réciproquement, si $\dimv(G)=2$, l'espace engendré par $u,v,w$ coïncide avec le plan $\Vect(u,v)$, donc $w\in\Vect(u,v)$.

\p{4} Si $w=au+bv$, alors
\[
v+w=au+(1+b)v.
\]

\p{5} Dans la base $(u,v)$, on a
\[
[u+v]_{(u,v)}=\binom{1}{1},
\qquad
[v+w]_{(u,v)}=\binom{a}{1+b}.
\]

\p{6} Les deux vecteurs $u+v$ et $v+w$ forment une base de $\Vect(u,v)$ si et seulement s'ils sont libres. Dans une base de plan, cela équivaut à dire que le déterminant de leurs coordonnées est non nul :
\[
\det\begin{pmatrix}
1 & a\\
1 & 1+b
\end{pmatrix}\neq 0.
\]
Cela prouve l'équivalence demandée.

\p{7} On calcule
\[
\det\begin{pmatrix}
1 & a\\
1 & 1+b
\end{pmatrix}=1+b-a.
\]
La condition se simplifie donc en
\[
1+b-a\neq 0.
\]

\p{8} Dans ce cas, comme $w\in\Vect(u,v)$, on a
\[
G=\Vect(u,v).
\]
Ainsi
\[
F=G
\iff F=\Vect(u,v)
\iff 1+b-a\neq 0.
\]
Donc, lorsque $w=au+bv$,
\[
F=G \iff a\neq 1+b.
\]

\p{9} Si $(u,v,w)$ est libre, alors c'est une famille libre de trois vecteurs. Donc
\[
\dimv(G)=3.
\]

\p{10} Par définition, $F$ est engendré par deux vecteurs, donc
\[
\dimv(F)\leq 2.
\]

\p{11} On a alors
\[
\dimv(F)\leq 2<3=\dimv(G),
\]
donc nécessairement
\[
F\neq G.
\]

\p{12} On a montré deux choses :
\begin{itemize}
    \item si $w\notin\Vect(u,v)$, alors $F\neq G$ ;
    \item si $w\in\Vect(u,v)$, alors $F=G$ équivaut à $1+b-a\neq 0$ lorsque $w=au+bv$.
\end{itemize}
Ainsi
\[
F=G
\iff \bigl(w\in\Vect(u,v)\text{ et }u+v,\,v+w\text{ forment une base de }\Vect(u,v)\bigr).
\]

\p{13} La condition finale s'écrit :
\[
F=G \iff \exists(a,b)\in\R^2,\quad w=au+bv\ \text{et}\ a\neq 1+b.
\]
Autrement dit,
\[
F=G \iff w=au+bv\text{ avec }1+b-a\neq 0.
\]

\p{14} Dans $\R^2$, avec
\[
u=(1,0),\qquad v=(0,1),\qquad w=(a,b),
\]
on a
\[
u+v=(1,1),
\qquad
v+w=(a,1+b).
\]
La condition
\[
\Vect(u+v,v+w)=\R^2
\]
équivaut à la liberté des deux vecteurs, donc au déterminant non nul :
\[
\det\begin{pmatrix}
1 & a\\
1 & 1+b
\end{pmatrix}=1+b-a\neq 0.
\]

\p{15} On retrouve exactement la même condition par calcul direct de déterminant dans la base canonique.

\p{16} Exemple où $F=G$ : choisir $a=0$, $b=0$, donc $w=0$. Alors
\[
1+b-a=1\neq 0,
\]
donc $F=G$.

Exemple où $F\neq G$ : choisir $a=1$, $b=0$. Alors
\[
1+b-a=0,
\]
donc $F\neq G$.



\end{document}
